\chapter{Wie funktioniert KI? Wie funktionieren LLMs?}\label{ch:K01_KI_Grundlagen}

\begin{myremark}
Dieses Kapitel führt in die Grundlagen der Künstlichen Intelligenz (KI) und insbesondere
in grosse Sprachmodelle (Large Language Models, LLMs) ein. Das Ziel ist ein intuitives,
schülernahes Verständnis – ohne tiefgreifende mathematische Vorkenntnisse.
\end{myremark}

% ─────────────────────────────────────────────────────────────────────────────
\section{Was ist Künstliche Intelligenz?}

\begin{mydefinition}
\textbf{Künstliche Intelligenz (KI)} bezeichnet Systeme, die Aufgaben ausführen, die
bisher typischerweise menschliche Intelligenz erforderten – zum Beispiel Sprache
verstehen, Bilder erkennen oder Entscheidungen treffen.

KI ist \emph{kein} magisches Denken: Im Kern handelt es sich immer um mathematische
Modelle, die aus Daten \enquote{lernen}.
\end{mydefinition}

\begin{myexample}
Bekannte KI-Systeme im Alltag:
\begin{itemize}
    \item \textbf{Empfehlungssysteme}: Netflix schlägt Filme vor, Spotify kuratiert Playlists.
    \item \textbf{Bilderkennungs-KI}: Das Smartphone entsperrt sich per Gesichtserkennung.
    \item \textbf{Sprach-KI}: Siri, Google Assistant oder Alexa beantworten gesprochene Fragen.
    \item \textbf{Große Sprachmodelle}: ChatGPT, Claude, Gemini führen Konversationen
          und generieren Texte.
\end{itemize}
\end{myexample}

\begin{myexercise}[– Brainstorming]
Sammeln Sie in 3 Minuten alle KI-Anwendungen, die Sie kennen oder täglich nutzen.
Vergleichen Sie anschliessend Ihre Liste mit einer Nachbarperson.
Welche KI-Anwendungen überraschten Sie?
\end{myexercise}

% ─────────────────────────────────────────────────────────────────────────────
\section{Maschinelles Lernen – Lernen aus Daten}

\begin{mydefinition}
\textbf{Maschinelles Lernen (ML)} ist ein Teilbereich der KI. Statt explizit
programmiert zu werden, \emph{lernt} ein ML-Modell Muster aus Beispieldaten
und kann danach Vorhersagen für neue, unbekannte Daten treffen.
\end{mydefinition}

\subsection{Ein einfaches Beispiel: Spam-Filter}

\begin{myexample}
Ein klassischer Spam-Filter wird nicht von Hand mit Regeln programmiert
(\enquote{Falls der Betreff das Wort \enquote{Gratis} enthält, ist es Spam}).
Stattdessen erhält das Modell tausende von E-Mails, die als \enquote{Spam} oder
\enquote{kein Spam} beschriftet sind, und \emph{lernt selbst} Muster daraus.

Dieses Prinzip nennt man \textbf{überwachtes Lernen} (supervised learning):
\begin{enumerate}
    \item \textbf{Trainingsdaten}: Viele beschriftete Beispiele (E-Mails mit Label).
    \item \textbf{Training}: Das Modell passt seine internen Parameter an, um
          möglichst viele Beispiele korrekt zu klassifizieren.
    \item \textbf{Vorhersage}: Das Modell entscheidet für neue E-Mails,
          ob es Spam ist oder nicht.
\end{enumerate}
\end{myexample}

\begin{myexercise}[– Entscheidungsbaum per Hand]
Stellen Sie sich vor, Sie möchten einem Kind erklären, wie es eine E-Mail als Spam
erkennen kann. Zeichnen Sie einen einfachen Entscheidungsbaum (Flussdiagramm) mit
mindestens drei Entscheidungsknoten und zwei möglichen Ausgaben (\enquote{Spam} /
\enquote{Kein Spam}).

Vergleichen Sie Ihre Lösung mit der Nachbarperson.
Was sind Vor- und Nachteile eines solchen handgefertigten Entscheidungsbaums gegenüber
einem maschinell gelernten Modell?
\end{myexercise}

% ─────────────────────────────────────────────────────────────────────────────
\section{Neuronale Netze – vom Gehirn inspiriert}

\begin{mydefinition}
\textbf{Künstliche neuronale Netze} sind lose vom menschlichen Gehirn inspiriert.
Sie bestehen aus \textbf{Knoten} (Neuronen) und \textbf{Verbindungen} (Gewichten).

\begin{itemize}
    \item Ein Neuron nimmt Eingaben entgegen, berechnet eine gewichtete Summe
          und gibt einen Wert weiter.
    \item Das Netz ist in \textbf{Schichten} aufgebaut:
          Eingabeschicht → versteckte Schichten → Ausgabeschicht.
    \item Die \textbf{Gewichte} werden durch Training angepasst, sodass das Netz
          die richtigen Ausgaben liefert.
\end{itemize}
\end{mydefinition}

\begin{myexample}[– Gewichte anpassen: ein Spielzeug-Neuron]
Stellen Sie sich ein einzelnes Neuron vor, das entscheidet, ob jemand ein Velo kaufen
wird. Es bekommt zwei Eingaben: \texttt{hat\_Geld} (0 oder 1) und
\texttt{mag\_Sport} (0 oder 1). Die Gewichte seien $w_1 = 0.8$ und $w_2 = 0.3$.

\[
    \text{Ausgabe} = w_1 \cdot \texttt{hat\_Geld} + w_2 \cdot \texttt{mag\_Sport}
\]

Falls die Ausgabe $> 0.5$ ist, sagt das Neuron \enquote{Ja, kauft ein Velo}.

Für eine Person mit Geld und Sportbegeisterung:
\[
    0.8 \cdot 1 + 0.3 \cdot 1 = 1.1 > 0.5 \quad \Rightarrow \quad \text{Ja}
\]
\end{myexample}

\begin{myexercise}[– Neuron berechnen]
Verwenden Sie das Spielzeug-Neuron aus dem Beispiel ($w_1 = 0.8$, $w_2 = 0.3$,
Schwelle $= 0.5$). Berechnen Sie die Ausgabe für die folgenden drei Personen und
entscheiden Sie, ob das Neuron \enquote{Ja} oder \enquote{Nein} sagt:

\begin{enumerate}
    \item Person A: hat Geld (\texttt{1}), mag Sport nicht (\texttt{0}).
    \item Person B: hat kein Geld (\texttt{0}), mag Sport (\texttt{1}).
    \item Person C: hat kein Geld (\texttt{0}), mag Sport nicht (\texttt{0}).
\end{enumerate}

Diskutieren Sie: Ist dieses Modell sinnvoll? Was würde sich ändern, wenn man
$w_2 = 0.6$ setzt?
\end{myexercise}

% ─────────────────────────────────────────────────────────────────────────────
\section{Wie funktionieren LLMs? – Grosse Sprachmodelle}

\begin{mydefinition}
\textbf{Large Language Models (LLMs)} sind sehr grosse neuronale Netze, die auf
riesigen Mengen von Textdaten trainiert wurden. Sie lernen, das \emph{nächste Wort}
(oder genauer: das nächste \textbf{Token}) in einem Text vorherzusagen.

Bekannte LLMs: GPT-4 (OpenAI), Claude (Anthropic), Gemini (Google), Llama (Meta).
\end{mydefinition}

\subsection{Token – die Bausteine eines LLMs}

\begin{mydefinition}
Ein \textbf{Token} ist die kleinste Einheit, mit der ein LLM arbeitet. Ein Token
entspricht ungefähr 3–4 Buchstaben oder einem kurzen Wort. Das Wort
\enquote{Informatik} könnte z.\,B. als ein einziges Token dargestellt werden,
während \enquote{unvorstellbar} in mehrere Tokens aufgeteilt werden kann.
\end{mydefinition}

\begin{myexample}[– Nächstes Token vorhersagen]
Das LLM sieht den Anfang eines Satzes:
\begin{center}
    \textit{Die Sonne scheint und der Himmel ist \ldots}
\end{center}
und muss das nächste wahrscheinlichste Token vorhersagen.
Mögliche Kandidaten und ihre geschätzten Wahrscheinlichkeiten:

\begin{center}
\begin{tabular}{ll}
    \toprule
    \textbf{Token} & \textbf{Wahrscheinlichkeit} \\
    \midrule
    blau           & 45\,\%                       \\
    wolkenlos      & 25\,\%                       \\
    klar           & 15\,\%                       \\
    strahlend      & 10\,\%                       \\
    grün           & 5\,\%                        \\
    \bottomrule
\end{tabular}
\end{center}

Das LLM wählt (meistens) das wahrscheinlichste Token – aber durch einen
\textbf{Temperatur}-Parameter kann auch ein weniger wahrscheinliches Token
gewählt werden, um kreativere Antworten zu erhalten.
\end{myexample}

\begin{myexercise}[– Nächstes Wort vorhersagen]
\textbf{Spielen Sie selbst LLM!} Lesen Sie die folgenden Satzanfänge und schreiben
Sie auf, welches Wort Sie als nächstes für am wahrscheinlichsten halten.
Vergleichen Sie Ihre Antworten in der Gruppe.

\begin{enumerate}
    \item \textit{Der kürzeste Weg zwischen zwei Punkten ist \ldots}
    \item \textit{In der Schweiz ist es verboten, \ldots}
    \item \textit{Das beliebteste soziale Netzwerk bei Jugendlichen ist \ldots}
    \item \textit{Wenn ich morgen krank bin, werde ich \ldots}
\end{enumerate}

Warum können Menschen dieses Spiel spielen? Was sagt das über LLMs aus?
\end{myexercise}

\subsection{Training von LLMs – wie lernen sie?}

\begin{myprocedure}[Training eines LLMs (vereinfacht)]
\begin{enumerate}
    \item \textbf{Daten sammeln}: Texte aus dem Internet, Bücher, Wikipedia
          (Billionen von Wörtern).
    \item \textbf{Vortraining} (self-supervised): Das Modell lernt, das nächste
          Token vorherzusagen. Fehler werden durch \textbf{Backpropagation}
          korrigiert – die Gewichte des Netzes werden angepasst.
    \item \textbf{Finetuning mit menschlichem Feedback (RLHF)}: Menschen bewerten
          Antworten. Das Modell lernt, hilfreichere, sicherere Antworten zu generieren.
    \item \textbf{Einsatz}: Das fertige Modell beantwortet Fragen in einem Chat.
\end{enumerate}
\end{myprocedure}

\begin{myremark}
Das Training grosser LLMs dauert \textbf{Monate} und kostet \textbf{Millionen von
Franken}. Es werden spezialisierte Chips (GPUs/TPUs) in riesigen Rechenzentren
verwendet. Einmal trainiert, kann das Modell allerdings von Millionen von Personen
gleichzeitig genutzt werden.
\end{myremark}

\begin{myexercise}[– Partnerübung: Was kann ein LLM (nicht)?]
Testen Sie folgende Aufgaben mit einem schulkonformen KI-Tool (z.\,B. Fobizz oder Brian)
und halten Sie das Ergebnis fest. Diskutieren Sie in Paaren:

\begin{enumerate}
    \item Bitten Sie die KI, ein Gedicht über den Zürichsee zu schreiben.
    \item Fragen Sie nach dem aktuellen Wetter in Ihrer Stadt.
    \item Bitten Sie die KI, eine schwierige Mathe-Gleichung zu lösen:
          $3x^2 - 5x + 2 = 0$.
    \item Fragen Sie: \enquote{Was ist Ihre Lieblingsfarbe?}
\end{enumerate}

\textbf{Reflexionsfragen}:
\begin{itemize}
    \item Bei welchen Aufgaben ist die KI stark / schwach?
    \item Was kann ein LLM prinzipiell \emph{nicht} wissen? Warum?
    \item Was bedeutet es, dass ein LLM \emph{kein} Bewusstsein hat?
\end{itemize}
\end{myexercise}

\subsection{Das Kontextfenster}

\begin{mydefinition}
Das \textbf{Kontextfenster} (context window) bestimmt, wie viel Text ein LLM auf
einmal \enquote{sehen} kann. Ältere Modelle hatten ein Fenster von wenigen tausend
Tokens. Moderne Modelle (z.\,B. Gemini 1.5 Pro) verarbeiten bis zu 1 Million Tokens –
das entspricht etwa 700 Buchseiten.
\end{mydefinition}

\begin{myremark}
Was \emph{ausserhalb} des Kontextfensters liegt, \enquote{vergisst} das Modell
vollständig. Ein LLM hat kein dauerhaftes Gedächtnis zwischen verschiedenen
Konversationen (ausser es wird explizit per Tool implementiert).
\end{myremark}

% ─────────────────────────────────────────────────────────────────────────────
\section{Wie unterscheidet sich KI von menschlichem Denken?}

\begin{myexercise}[– Gruppenvergleich]
Bearbeiten Sie in Vierergruppen folgende Aufgabe: Vergleichen Sie menschliches Denken
und das \enquote{Denken} eines LLMs anhand der folgenden Kriterien.
Füllen Sie die Tabelle aus:

\begin{center}
\begin{tabular}{p{4cm}p{4.5cm}p{4.5cm}}
    \toprule
    \textbf{Kriterium}           & \textbf{Mensch} & \textbf{LLM} \\
    \midrule
    Lernen                       & &                              \\
    Gedächtnis                   & &                              \\
    Kreativität                  & &                              \\
    Fehler machen                & &                              \\
    Emotionen                    & &                              \\
    Selbstbewusstsein            & &                              \\
    Aktualität der Informationen & &                              \\
    \bottomrule
\end{tabular}
\end{center}

Präsentieren Sie Ihre Ergebnisse der Klasse in 2 Minuten.
\end{myexercise}

\begin{mychallenge}[Erkunde einen Tokenizer]
Besuchen Sie den öffentlich zugänglichen Tokenizer von OpenAI:
\begin{center}
    \url{https://platform.openai.com/tokenizer}
\end{center}
Geben Sie verschiedene Texte ein (Deutsch, Englisch, Emojis, Code) und beobachten Sie,
wie der Text in Tokens aufgeteilt wird.

\begin{enumerate}
    \item Welches Wort braucht die meisten Tokens?
    \item Warum braucht englischer Text im Allgemeinen weniger Tokens als
          gleichwertiger deutscher Text?
    \item Was passiert mit Emojis?
\end{enumerate}
\end{mychallenge}
